% GENERATED FILE. Edit canonical lessons, then run npm run generate:books.
% canonical-source-hash: fnv1a64:ce89d2ac207d7927
% canonical-lessons: LA-C338-legio, LA-C338-cohors, LA-C338-centurio, LA-C338-signifer, LA-C338-tribunus, LA-R338-review-331-338

\chapter{Legiō}
\label{ch:la-legio}
\hlchaptermodality{338}

\begin{chapteropening}
\textbf{By the end of this chapter:} \emph{I can name the parts of a Roman army: the legion, the cohort, the centurion, the standard-bearer and the tribune.}

Complete the last lesson of chapter 338: I can name the parts of a Roman army: the legion, the cohort, the centurion, the standard-bearer and the tribune.
\end{chapteropening}

\section[legiō, legiōnis]{legiō, legiōnis — a legion}

\label{lesson:LA-C338-legio}



\begin{quote}
Before the new one: say the Latin for a contest, then the Latin for applause.
\end{quote}

\subsection*{You'll want to know: legiō, legiōnis}
\textbf{legiō, legiōnis} — \textquotedblleft{}a legion\textquotedblright{}.

\begin{quote}
Decima legiō fortissima erat.
\end{quote}

\begin{quote}
\emph{The tenth legion was the bravest.}
\end{quote}

Feminine, 3rd declension, from \emph{legō}, 'to choose': a picked levy, about five thousand men. English \emph{legion}.

\begin{cousinweb}
\textbf{legiō} is built on \emph{legō}, \textquotedblleft{}to gather, to choose\textquotedblright{} (Lewis \& Short): a body of chosen soldiers. English \emph{legion} comes from it through French.
\end{cousinweb}

\begin{grammarlens}[title={What you've built}]
One more word for telling, explaining and describing at B1.
\end{grammarlens}

\subsection*{Your turn}
\begin{itemize}
\raggedright
  \item \emph{Say it:} \emph{legiō, legiōnis}
  \item \emph{Say it:} \emph{legiō, legiōnis}, once more
  \item \emph{Say it:} a sentence of your own with \emph{legiō, legiōnis}
  \item \emph{Recall:} say the Latin for a contest, then the Latin for applause, then say \emph{legiō, legiōnis} again
\end{itemize}

\begin{tcolorbox}[breakable,colback=teal!4,colframe=teal!35!black,title={Before you move on}]
What does legiō, legiōnis mean? (\textquotedblleft{}A legion\textquotedblright{}.) Say it once more.
\end{tcolorbox}
\section[cohors, cohortis]{cohors, cohortis — a cohort; an enclosure}

\label{lesson:LA-C338-cohors}



\begin{quote}
Before the new one: say the Latin for applause, then the Latin for a legion.
\end{quote}

\subsection*{You'll want to know: cohors, cohortis}
\textbf{cohors, cohortis} — \textquotedblleft{}a cohort; an enclosure\textquotedblright{}.

\begin{quote}
Cohors ad portam castrōrum stābat.
\end{quote}

\begin{quote}
\emph{A cohort was standing at the gate of the camp.}
\end{quote}

Feminine, 3rd declension. A legion had ten cohorts. The same word meant a farmyard, which gives English \emph{court}.

\begin{cousinweb}
\textbf{cohors} is inherited, and its earlier history is not settled. It first meant \textquotedblleft{}an enclosure, a farmyard\textquotedblright{} (EtymDB, Lewis \& Short), and then the body of soldiers drawn up together. English \emph{court} comes from it through French, and \emph{cohort} directly.
\end{cousinweb}

\begin{grammarlens}[title={What you've built}]
One more word for telling, explaining and describing at B1.
\end{grammarlens}

\subsection*{Your turn}
\begin{itemize}
\raggedright
  \item \emph{Say it:} \emph{cohors, cohortis}
  \item \emph{Say it:} \emph{cohors, cohortis}, once more
  \item \emph{Say it:} a sentence of your own with \emph{cohors, cohortis}
  \item \emph{Recall:} say the Latin for applause, then the Latin for a legion, then say \emph{cohors, cohortis} again
\end{itemize}

\begin{tcolorbox}[breakable,colback=teal!4,colframe=teal!35!black,title={Before you move on}]
What does cohors, cohortis mean? (\textquotedblleft{}A cohort; an enclosure\textquotedblright{}.) Say it once more.
\end{tcolorbox}
\section[centuriō, centuriōnis]{centuriō, centuriōnis — a centurion}

\label{lesson:LA-C338-centurio}



\begin{quote}
Before the new one: say the Latin for a legion, then the Latin for a cohort.
\end{quote}

\subsection*{You'll want to know: centuriō, centuriōnis}
\textbf{centuriō, centuriōnis} — \textquotedblleft{}a centurion\textquotedblright{}.

\begin{quote}
Centuriō mīlitēs suōs ad mūrum dūxit.
\end{quote}

\begin{quote}
\emph{The centurion led his soldiers to the wall.}
\end{quote}

Masculine, 3rd declension, from \emph{centuria}, a company of a hundred.

\begin{cousinweb}
\textbf{centuriō} is built on \emph{centuria}, \textquotedblleft{}a division of a hundred\textquotedblright{}, from \emph{centum}, \textquotedblleft{}hundred\textquotedblright{} (EtymDB).
\end{cousinweb}

\begin{grammarlens}[title={What you've built}]
One more word for telling, explaining and describing at B1.
\end{grammarlens}

\subsection*{Your turn}
\begin{itemize}
\raggedright
  \item \emph{Say it:} \emph{centuriō, centuriōnis}
  \item \emph{Say it:} \emph{centuriō, centuriōnis}, once more
  \item \emph{Say it:} a sentence of your own with \emph{centuriō, centuriōnis}
  \item \emph{Recall:} say the Latin for a legion, then the Latin for a cohort, then say \emph{centuriō, centuriōnis} again
\end{itemize}

\begin{tcolorbox}[breakable,colback=teal!4,colframe=teal!35!black,title={Before you move on}]
What does centuriō, centuriōnis mean? (\textquotedblleft{}A centurion\textquotedblright{}.) Say it once more.
\end{tcolorbox}
\section[signifer, signiferī]{signifer, signiferī — a standard-bearer}

\label{lesson:LA-C338-signifer}



\begin{quote}
Before the new one: say the Latin for a cohort, then the Latin for a centurion.
\end{quote}

\subsection*{You'll want to know: signifer, signiferī}
\textbf{signifer, signiferī} — \textquotedblleft{}a standard-bearer\textquotedblright{}.

\begin{quote}
Signifer signum in hostēs tulit.
\end{quote}

\begin{quote}
\emph{The standard-bearer carried the standard against the enemy.}
\end{quote}

Masculine: \emph{signum} (taught) + \emph{ferō}. Losing the standard was a deep disgrace.

\begin{cousinweb}
\textbf{signifer} joins \emph{signum}, \textquotedblleft{}sign, standard\textquotedblright{}, and \emph{-fer}, \textquotedblleft{}carrying\textquotedblright{}, from \emph{ferō} (EtymDB).
\end{cousinweb}

\begin{grammarlens}[title={What you've built}]
One more word for telling, explaining and describing at B1.
\end{grammarlens}

\subsection*{Your turn}
\begin{itemize}
\raggedright
  \item \emph{Say it:} \emph{signifer, signiferī}
  \item \emph{Say it:} \emph{signifer, signiferī}, once more
  \item \emph{Say it:} a sentence of your own with \emph{signifer, signiferī}
  \item \emph{Recall:} say the Latin for a cohort, then the Latin for a centurion, then say \emph{signifer, signiferī} again
\end{itemize}

\begin{tcolorbox}[breakable,colback=teal!4,colframe=teal!35!black,title={Before you move on}]
What does signifer, signiferī mean? (\textquotedblleft{}A standard-bearer\textquotedblright{}.) Say it once more.
\end{tcolorbox}
\section[tribūnus, tribūnī]{tribūnus, tribūnī — a tribune}

\label{lesson:LA-C338-tribunus}



\begin{quote}
Before the new one: say the Latin for a centurion, then the Latin for a standard-bearer.
\end{quote}

\subsection*{You'll want to know: tribūnus, tribūnī}
\textbf{tribūnus, tribūnī} — \textquotedblleft{}a tribune\textquotedblright{}.

\begin{quote}
Tribūnus mīlitum ad castra advēnit.
\end{quote}

\begin{quote}
\emph{The military tribune arrived at the camp.}
\end{quote}

Masculine. \emph{Tribūnus mīlitum}: an officer of a legion; the people's tribunes protected ordinary citizens.

\begin{cousinweb}
\textbf{tribūnus} is built on \emph{tribus}, \textquotedblleft{}tribe\textquotedblright{} (EtymDB): first the head of a tribe. English \emph{tribune} comes from it.
\end{cousinweb}

\begin{grammarlens}[title={What you've built}]
One more word for telling, explaining and describing at B1.
\end{grammarlens}

\subsection*{Your turn}
\begin{itemize}
\raggedright
  \item \emph{Say it:} \emph{tribūnus, tribūnī}
  \item \emph{Say it:} \emph{tribūnus, tribūnī}, once more
  \item \emph{Say it:} a sentence of your own with \emph{tribūnus, tribūnī}
  \item \emph{Recall:} say the Latin for a centurion, then the Latin for a standard-bearer, then say \emph{tribūnus, tribūnī} again
\end{itemize}

\begin{tcolorbox}[breakable,colback=teal!4,colframe=teal!35!black,title={Before you move on}]
What does tribūnus, tribūnī mean? (\textquotedblleft{}A tribune\textquotedblright{}.) Say it once more.
\end{tcolorbox}
\section[(dialogue)]{Words from chapters 331–338 — a review}

\label{lesson:LA-R338-review-331-338}



\begin{quote}
Say the Latin for grief and a dirge.
\end{quote}

\subsection*{Your turn}
Say these aloud:

\begin{itemize}
\raggedright
  \item grief, a dirge, a speech of praise, a memorial, offerings to the dead
  \item an heir, an inheritance, a will, to seal, to leave behind
  \item festive, yearly, a procession, a feast, to feast
  \item a priest, a sacrificial animal, to sacrifice, to vow, a vow
  \item a charioteer, a four-horse chariot, a turning post, a racing team, of the circus
  \item a gladiator, a trainer of gladiators, a public show, a hunt, a beast
  \item the palm of victory, a prize, to compete, a contest, applause
  \item a legion, a cohort, a centurion, a standard-bearer, a tribune
\end{itemize}

\begin{tcolorbox}[breakable,colback=teal!4,colframe=teal!35!black,title={Before you move on}]
Grief? (\textbf{lūctus, lūctūs}.) A dirge? (\textbf{nēnia, nēniae}.) A speech of praise? (\textbf{laudātiō, laudātiōnis}.) A memorial? (\textbf{monumentum, monumentī}.) Offerings to the dead? (\textbf{īnferiae, īnferiārum}.) An heir? (\textbf{hērēs, hērēdis}.) An inheritance? (\textbf{hērēditās, hērēditātis}.) A will? (\textbf{testāmentum, testāmentī}.) To seal? (\textbf{obsignō, obsignāre}.) To leave behind? (\textbf{relinquō, relinquere}.) Festive? (\textbf{fēstus, fēsta, fēstum}.) Yearly? (\textbf{sollemnis, sollemne}.) A procession? (\textbf{pompa, pompae}.) A feast? (\textbf{epulae, epulārum}.) To feast? (\textbf{epulor, epulārī}.) A priest? (\textbf{sacerdōs, sacerdōtis}.) A sacrificial animal? (\textbf{victima, victimae}.) To sacrifice? (\textbf{immolō, immolāre}.) To vow? (\textbf{voveō, vovēre}.) A vow? (\textbf{vōtum, vōtī}.) A charioteer? (\textbf{aurīga, aurīgae}.) A four-horse chariot? (\textbf{quadrīgae, quadrīgārum}.) A turning post? (\textbf{mēta, mētae}.) A racing team? (\textbf{factiō, factiōnis}.) Of the circus? (\textbf{circēnsis, circēnse}.) A gladiator? (\textbf{gladiātor, gladiātōris}.) A trainer of gladiators? (\textbf{lanista, lanistae}.) A public show? (\textbf{mūnus, mūneris}.) A hunt? (\textbf{vēnātiō, vēnātiōnis}.) A beast? (\textbf{bēstia, bēstiae}.) The palm of victory? (\textbf{palma, palmae}.) A prize? (\textbf{praemium, praemiī}.) To compete? (\textbf{certō, certāre}.) A contest? (\textbf{certāmen, certāminis}.) Applause? (\textbf{plausus, plausūs}.) A legion? (\textbf{legiō, legiōnis}.) A cohort? (\textbf{cohors, cohortis}.) A centurion? (\textbf{centuriō, centuriōnis}.) A standard-bearer? (\textbf{signifer, signiferī}.) A tribune? (\textbf{tribūnus, tribūnī}.)
\end{tcolorbox}
